\uuid{lOW2}
\titre{Transformations et propriétés de la loi log-normale}

\niveau{L3}
\module{Modélisation de l'aléatoire}
\chapitre{Transformations de variables aléatoires}
\sousChapitre{Loi log-normale}
\theme{Inverse, produit et espérance d'une variable log-normale}
\auteur{Erwan Hillion}
\datecreate{2026-10-01}
\organisation{ESM3/ISD}
\difficulte{3}

\contenu{

\texte{
La loi log-normale de paramètres $\mu\in\mathbb R$ et $\sigma>0$, notée $LN(\mu,\sigma)$, est la loi de densité
$$
f(x)=\frac{1}{x\sigma\sqrt{2\pi}}
\exp\left(-\frac{(\ln x-\mu)^2}{2\sigma^2}\right)
\mathbf 1_{]0,+\infty[}(x).
$$
Dans cet exercice, le second paramètre de $\mathcal N(\mu,\sigma)$ désigne l'écart-type.
}

\begin{enumerate}

\item \question{Soit $X$ une variable aléatoire à valeurs strictement positives. Montrer que $Y=\ln X$ suit la loi normale $\mathcal N(\mu,\sigma)$ si et seulement si $X$ suit la loi $LN(\mu,\sigma)$.}
\indication{}
\reponse{
Supposons d'abord que $X\sim LN(\mu,\sigma)$ et posons $Y=\ln X$. Le changement de variable $x=e^y$ donne $dx=e^y\,dy$. Ainsi,
\begin{align*}
f_Y(y)
&=f_X(e^y)e^y\\
&=\frac{1}{e^y\sigma\sqrt{2\pi}}
\exp\left(-\frac{(y-\mu)^2}{2\sigma^2}\right)e^y\\
&=\frac{1}{\sigma\sqrt{2\pi}}
\exp\left(-\frac{(y-\mu)^2}{2\sigma^2}\right).
\end{align*}
Donc $Y\sim\mathcal N(\mu,\sigma)$.

Réciproquement, supposons $Y\sim\mathcal N(\mu,\sigma)$ et posons $X=e^Y$. Pour $x>0$, le changement de variable $y=\ln x$ donne $dy=dx/x$, donc
$$
f_X(x)=\frac1x f_Y(\ln x)
=\frac{1}{x\sigma\sqrt{2\pi}}
\exp\left(-\frac{(\ln x-\mu)^2}{2\sigma^2}\right).
$$
Ainsi,
$$
\boxed{Y=\ln X\sim\mathcal N(\mu,\sigma)
\iff X\sim LN(\mu,\sigma).}
$$
}

\item \question{Montrer que si $X\sim LN(\mu,\sigma)$, alors $1/X$ suit une loi log-normale dont on précisera les paramètres.}
\indication{}
\reponse{
On a
$$
\ln\left(\frac1X\right)=-\ln X.
$$
Or, si $X\sim LN(\mu,\sigma)$, alors $\ln X\sim\mathcal N(\mu,\sigma)$. Par conséquent,
$$
-\ln X\sim\mathcal N(-\mu,\sigma).
$$
D'après la caractérisation précédente,
$$
\boxed{\frac1X\sim LN(-\mu,\sigma).}
$$
}

\item \question{Soit $(X_i)_{1\leq i\leq n}$ une famille de variables aléatoires mutuellement indépendantes telle que $X_i\sim LN(\mu_i,\sigma_i)$. Déterminer la loi de
$$
Z=\prod_{i=1}^nX_i.
$$}
\indication{}
\reponse{
Posons $Y_i=\ln X_i$. Alors les variables $Y_i$ sont indépendantes et
$$
Y_i\sim\mathcal N(\mu_i,\sigma_i).
$$
De plus,
$$
\ln Z=\sum_{i=1}^n\ln X_i=\sum_{i=1}^nY_i.
$$
Une somme de variables normales indépendantes est normale. On a donc
$$
\ln Z\sim\mathcal N\left(
\sum_{i=1}^n\mu_i,
\sqrt{\sum_{i=1}^n\sigma_i^2}
\right).
$$
Par conséquent,
$$
\boxed{
Z\sim LN\left(
\sum_{i=1}^n\mu_i,
\sqrt{\sum_{i=1}^n\sigma_i^2}
\right).
}
$$
}

\item \question{On suppose que $X\sim LN(5,1{,}5)$. Calculer $\mathbb P(50\leq X\leq400)$.}
\indication{}
\reponse{
La variable $Y=\ln X$ suit la loi $\mathcal N(5,1{,}5)$. Ainsi,
\begin{align*}
\mathbb P(50\leq X\leq400)
&=\mathbb P(\ln50\leq Y\leq\ln400)\\
&=\Phi\left(\frac{\ln400-5}{1{,}5}\right)
-\Phi\left(\frac{\ln50-5}{1{,}5}\right),
\end{align*}
où $\Phi$ désigne la fonction de répartition de la loi normale centrée réduite.
Numériquement,
$$
\frac{\ln50-5}{1{,}5}\simeq-0{,}725,
\qquad
\frac{\ln400-5}{1{,}5}\simeq0{,}661.
$$
Donc
$$
\boxed{
\mathbb P(50\leq X\leq400)
\simeq\Phi(0{,}661)-\Phi(-0{,}725)
\simeq0{,}512.
}
$$
}

\item \question{Soit $X\sim LN(0,1)$. Montrer que
$$
\mathbb E[X]
=\frac{1}{\sqrt{2\pi}}\int_{\mathbb R}e^{y-y^2/2}\,dy
=e^{1/2}.
$$}
\indication{}
\reponse{
Si $Y=\ln X$, alors $Y\sim\mathcal N(0,1)$ et $X=e^Y$. Par conséquent,
$$
\mathbb E[X]
=\mathbb E[e^Y]
=\frac{1}{\sqrt{2\pi}}
\int_{\mathbb R}e^y e^{-y^2/2}\,dy.
$$
On complète le carré :
$$
y-\frac{y^2}{2}
=-\frac{(y-1)^2}{2}+\frac12.
$$
Ainsi,
\begin{align*}
\mathbb E[X]
&=\frac{e^{1/2}}{\sqrt{2\pi}}
\int_{\mathbb R}e^{-(y-1)^2/2}\,dy\\
&=e^{1/2},
\end{align*}
car l'intégrale restante est celle d'une densité normale de moyenne $1$ et d'écart-type $1$.
Donc
$$
\boxed{\mathbb E[X]=e^{1/2}.}
$$
}

\end{enumerate}

}
